\section{二分图最大权匹配（KM）}

KM 求两侧各 $n$ 点的二分图最大权（完美）匹配，要求存在由实边组成的完美匹配。权矩阵 \texttt{w} 按下标 $1..n$ 填：\verb|addEdge(u,v,c)| 在 $u$ 与 $v$ 间建权 $c$ 的边，重边取较大者；未建边保持初始化极小值 $-\infty$（约 $-4\times 10^{18}$），故负权实边也能正常加入，最优完美匹配不会用到极小值边。两侧点数不等时，先给较小一侧补上数量差个权 $0$ 的虚拟点（相应 $0$ 边照常 \verb|addEdge|），补到同阶后结果即为两侧各自的“最大权匹配”。求最小权时把所有权取负后调用。

\verb|work()| 返回最大权和；结束后 \verb|match[j]| 为右侧 $j$ 匹配到的左侧点。

复杂度 $\mathcal{O}(n^3)$。

\begin{lstlisting}
using i64 = long long;
struct KM {
    int n;
    std::vector<std::vector<i64>> w;
    std::vector<i64> lx, ly, slk;
    std::vector<int> match, pre;
    std::vector<char> vis;
    KM(int n) : n(n), w(n + 1, std::vector<i64>(n + 1, -(i64)4e18)),
                lx(n + 1), ly(n + 1), slk(n + 1),
                match(n + 1), pre(n + 1), vis(n + 1) {}
    void addEdge(int u, int v, i64 c){ w[u][v] = std::max(w[u][v], c); }
    void augment(int root){
        const i64 INF = (i64)4e18;
        std::fill(vis.begin(), vis.end(), 0);
        std::fill(slk.begin(), slk.end(), INF);
        int y = 0, yy = 0;
        match[0] = root;
        do {
            int a = match[y];
            i64 d = INF;
            vis[y] = 1;
            for (int b = 1; b <= n; b++){
                if (vis[b]) continue;
                i64 g = lx[a] + ly[b] - w[a][b];
                if (g < slk[b]){ slk[b] = g; pre[b] = y; }
                if (slk[b] < d){ d = slk[b]; yy = b; }
            }
            for (int b = 0; b <= n; b++){
                if (vis[b]){ lx[match[b]] -= d; ly[b] += d; }
                else slk[b] -= d;
            }
            y = yy;
        } while (match[y]);
        while (y){
            match[y] = match[pre[y]];
            y = pre[y];
        }
    }
    i64 work(){
        const i64 INF = (i64)4e18;
        for (int i = 1; i <= n; i++){
            lx[i] = -INF;
            for (int j = 1; j <= n; j++) lx[i] = std::max(lx[i], w[i][j]);
            ly[i] = 0;
        }
        std::fill(match.begin(), match.end(), 0);
        for (int i = 1; i <= n; i++) augment(i);
        i64 ans = 0;
        for (int j = 1; j <= n; j++) ans += w[match[j]][j];
        return ans;
    }
};
\end{lstlisting}
